% PowerPoint の標準と同じ 16:9
\documentclass[aspectratio=169]{beamer}
\usepackage{tikz}

% 数式はセリフ体（Computer Modern系）に統一し、
% Beamer既定のサンセリフ数式を使わない
\usefonttheme{professionalfonts}
\usefonttheme[onlymath]{serif}

\usepackage{luatexja}
\usepackage{luatexja-fontspec}
\usepackage{graphicx}
\usepackage{bm}
% 章フォルダの tex/ からコンパイルすることを想定
\graphicspath{{../images/}{images/}}

% ベクトル・パラメータ表記を共通化
\newcommand{\vx}{\bm{x}}
\newcommand{\vu}{\bm{u}}
\newcommand{\veps}{\bm{\epsilon}}
\newcommand{\vw}{\bm{w}}

% ---- ブロックの設定（block / alertblock / exampleblock）----
\definecolor{main}{RGB}{30,70,140}
\definecolor{alertcol}{RGB}{200,80,40}
\definecolor{examplecol}{RGB}{30,120,60}
\setbeamertemplate{blocks}[rounded][shadow=false]

\setbeamercolor{block title}{fg=white,bg=main}
\setbeamercolor{block body}{fg=black,bg=main!6}
\setbeamercolor{block title alerted}{fg=white,bg=alertcol}
\setbeamercolor{block body alerted}{fg=black,bg=alertcol!6}
\setbeamercolor{block title example}{fg=white,bg=examplecol}
\setbeamercolor{block body example}{fg=black,bg=examplecol!6}

\setbeamerfont{block title}{series=\bfseries}
\setbeamerfont{block title alerted}{series=\bfseries}
\setbeamerfont{block title example}{series=\bfseries}

\setbeamertemplate{itemize items}[circle]
\setbeamertemplate{navigation symbols}{}
\setbeamercolor{page number in head/foot}{fg=black!60}
\setbeamertemplate{footline}[frame number]

\DeclareMathOperator*{\argmax}{arg\,max}
\DeclareMathOperator*{\argmin}{arg\,min}

\author{渡邉 有志}
\date{\today}
